\chapter{Literature Review}

\section{Review and Comparison of Previous Work}
The landscape of Security Operations Centers (SOC) and Security Information and Event Management (SIEM) solutions has undergone significant structural evolution over the past two decades. The prevailing paradigm in enterprise cybersecurity has consistently advocated for deploying highly centralised, monolithic data collection frameworks explicitly designed to aggregate, index, and query vast swathes of telemetry generated across disparate servers, endpoint nodes, and network appliances (firewalls, routers, switches).

\subsection{Intrusion Detection Systems (IDS)}
The foundational science of anomaly detection originates in the architecture of Intrusion Detection Systems (IDS), historically categorized into two functionally opposed implementations: Network-based Intrusion Detection Systems (NIDS) and Host-based Intrusion Detection Systems (HIDS). 

NIDS platforms, typified by open-source solutions like Snort or Suricata, analyze network traffic transiting a physical switch via port mirroring or TAP devices. While incredibly effective at identifying malicious packet signatures, NIDS visibility ends at the encryption layer. As modern internet traffic is overwhelmingly encrypted (HTTPS, TLS 1.3), NIDS are largely blinded to the actual payloads unless computationally expensive SSL decryption appliances are deployed.

Conversely, HIDS environments function natively upon the operating system of the target endpoint. These agents bypass network encryption because they process data at the application layer, scrutinizing internal artifacts such as memory allocation, process executions, file integrity, and authentication attempts. Projects like OSSEC pioneered this domain. HIDS provide unparalleled depth of visibility; however, distributing, updating, and coordinating thousands of independent host agents present a colossal logistical challenge for network administrators. The HomeSOC project inherits its conceptual lineage from the HIDS methodology, prioritizing endpoint depth over shallow network breadth.

\subsection{Enterprise Security Information and Event Management (SIEM)}
The integration of centralized log indexing with HIDS alerts birthed the modern SIEM. Prominent examples of conventional SIEM architectures include commercial industry-standard platforms such as Splunk, IBM QRadar, and open-source equivalents like the ELK stack (Elasticsearch, Logstash, Kibana).

Extensive academic studies and industry deployments have historically documented the implementation of the ELK stack as the preeminent methodology for parsing arbitrary log formats and creating complex geographical and statistical dashboards. In these models, lightweight network shippers (most notably components of the Beats family, such as Filebeat or Metricbeat) continuously transmit raw event logs in near-real-time to a central indexing engine (Elasticsearch). The data is tokenised, inverted, and stored across distributed shards, making it instantaneously searchable.

While this architecture is highly customizable and horizontally scalable, enterprise-grade systems inherently require exorbitant computational resources to sustain their Java Virtual Machine (JVM) dependencies. Furthermore, they demand intricate configuration pipelines utilizing complex Grok filtering or regular expressions (REGEX) to coerce unstructured syslog inputs into structured JSON. The infrastructure capital expenditure to maintain a fault-tolerant Elasticsearch cluster—characterised by its extensive RAM consumption—renders it largely unsuitable and financially unviable for low-power environments, small offices, or hobbyist local networks.

\subsection{Wazuh: Bridging SIEM and HIDS}
Recognizing the limitations and disjointed nature of standalone HIDS and raw SIEM loggers, Wazuh emerged as a viable, open-source centralized security monitoring platform. Originating as a fork of the aging OSSEC project, Wazuh successfully merges SIEM capabilities with agent-based compliance monitoring, delivering robust endpoint detection, hardware vulnerability scanning, and File Integrity Monitoring (FIM).

However, despite its open-source license, the Wazuh managerial architecture still fundamentally leverages an intricate stack of dependencies—reliant on Filebeat for transport and an Elasticsearch indexing cluster (or recently, the OpenSearch equivalent) for archival indexing. Requiring substantial system administration expertise to deploy the multi-node infrastructure, interpret the XML-based rule syntax, and navigate the Kibana UI plugins, the inherent complexity presents a monumental barrier to entry. This friction consistently prevents non-expert users, small business operators, or residential network owners from attempting a deployment, leaving their critical digital assets unmonitored and vulnerable.

\begin{figure}[h]
    \centering
    % Placeholder for Wazuh/ELK Screenshot
    \fbox{
        \begin{minipage}{0.8\textwidth}
            \vspace{2cm}
            \begin{center}
                \textit{[PLACEHOLDER: Insert Screenshot of complex Wazuh/Kibana UI Dashboard here]} \\
                \vspace{0.5cm}
                File: \texttt{wazuh\_dashboard\_example.png}
            \end{center}
            \vspace{2cm}
        \end{minipage}
    }
    \caption{Example of the dense, expert-oriented UI design typically found in enterprise SIEM stacks.}
    \label{fig:siem_ui}
\end{figure}

\section{Highlight and Comparison of Proposed Study against Previous Work}
While traditional SIEM architectures prioritise total visibility across global, distributed enterprise assets regardless of processing cost, the proposed ``HomeSOC'' study deliberately shifts the technological paradigm toward extreme efficiency, modularity, and accessibility. 

The most pronounced distinction between HomeSOC and preceding systems lies in the literal architectural footprint and baseline requirement of the infrastructure. Unlike the ELK stack or Wazuh, which demand multi-gigabyte memory allocations merely to sustain their core database listener processes, the proposed system operates entirely on native, lightweight Python libraries and a highly engineered local SQLite backend. By capitalising on SQLite's robust, zero-configuration serverless architecture, and manually enabling Write-Ahead Logging (WAL) concurrency, the database processing engine in HomeSOC achieves robust API ingestion parity without the overhead of maintaining a separate, daemonised database service (such as MySQL or PostgreSQL).

Table 2-1 synthesizes the comparative resource overhead of HomeSOC against common industry standards.

\begin{table}[h]
\centering
\small
\begin{tabularx}{\textwidth}{|X|c|c|X|}
\hline
\textbf{Architectural Metric} & \textbf{Elastic ELK Stack} & \textbf{Wazuh Platform} & \textbf{HomeSOC System} \\ \hline
Minimum System RAM & 4GB - 8GB & 2GB - 4GB & 256MB - 512MB \\ \hline
Overall Architecture & Distributed Monolith & Manager-Agent Hybrid & Lightweight Client-Server \\ \hline
Database Core & Elasticsearch Engine & Indexer / MySQL & SQLite (WAL Pragmas) \\ \hline
Deployment Complexity & Extremely High & Medium-High & Low (Script-based Exec) \\ \hline
Primary Audience & SOC Analysts & IT Administrators & Non-Experts / SME \\ \hline
\end{tabularx}
\caption{Comparative analytical breakdown of HomeSOC physical/virtual resource requirements against prevailing enterprise SOC frameworks.}
\end{table}

\subsection{Heuristic Processing vs. Indexing Search}
Furthermore, traditional systems pass the burden of threat classification onto complex indexing backends. They ingest everything natively and rely on analysts to write complex search queries (e.g., KQL or Lucene syntax) against the massive dataset. The proposed study seamlessly integrates a custom detection script (`detection.py`) directly within the backend Flask application pipeline. 

Instead of caching gigabytes of routine "normal" behavior data, the HomeSOC host agent focuses aggressively on high-fidelity, specific telemetry vectors—such as targeted authentication failures, critical specific process execution (`cmd.exe`, `powershell`, `nmap`), and predefined mathematically-based threshold parameters. The volume of data transmitted over the network is drastically minimized via this filter. This architectural choice severely constrains network latency and produces entirely negligible impact on host processing cycles. 

\subsection{Democratizing The UI Experience}
Ultimately, previous works primarily focused on scaling complex big-data pipelines for trained cybersecurity analysts operating in Tier-1 and Tier-2 data centers. In stark contrast, this study prioritises sheer aesthetic usability, responsive design, and rapid triage for absolute laypersons. 

The HomeSOC implementation curates a utilitarian, custom dark-mode dashboard (devoid of heavy frameworks like React or Angular) that dynamically distils raw mathematical JSON events into concise, actionable layman remediation steps. These summaries are accompanied by highly visible severity toasts and animated timelines. By fundamentally translating cryptography and cyber events into plain English, the approach democratizes cybersecurity monitoring. The system cleanly fulfills the critical need for a lightweight, strictly functional, and easy-to-deploy local command centre that can be realistically sustained on cloud infrastructure like PythonAnywhere.

\section{Expanded Theoretical Framework: System and Intrusion Methodologies}
Security information and event management (SIEM) is a field within computer security that combines security information management (SIM) and security event management (SEM) to enable real-time analysis of security alerts generated by applications and network hardware. SIEM systems are central to security operations centers (SOCs), where they are employed to detect, investigate, and respond to security incidents. SIEM technology collects and aggregates data from various systems, allowing organizations to meet compliance requirements while safeguarding against threats. NIST's definition for a SIEM tool is an application that provides the ability to gather security data from information system components and present that data as actionable information via a single interface.
SIEM tools can be implemented as software, hardware, or managed services. SIEM systems log security events and generate reports to meet regulatory frameworks such as the Health Insurance Portability and Accountability Act (HIPAA) and the Payment Card Industry Data Security Standard (PCI DSS). The integration of SIM and SEM within SIEM provides organizations with a centralized approach for monitoring security events and responding to threats in real-time.
First introduced by Gartner analysts Mark Nicolett and Amrit Williams in 2005, the term SIEM has evolved to incorporate advanced features such as threat intelligence and behavioral analytics, which allow SIEM solutions to manage complex cybersecurity threats, including zero-day vulnerabilities and polymorphic malware.
In recent years, SIEM has become increasingly incorporated into national cybersecurity initiatives. For instance, Executive Order 14028 signed in 2021 by U.S. President Joseph Biden mandates the use of SIEM technologies to improve incident detection and reporting in federal systems. Compliance with these mandates is further reinforced by frameworks such as NIST SP 800-92, which outlines best practices for managing computer security logs.
Modern SIEM platforms are aggregating and normalizing data not only from various Information Technology (IT) sources, but from production and manufacturing Operational Technology (OT) environments as well.


== History ==
Initially, system logging was primarily used for troubleshooting and debugging. However, as operating systems and networks have grown more complex, so has the generation of system logs. The monitoring of system logs has also become increasingly common due to the rise of sophisticated cyberattacks and the need for compliance with regulatory frameworks, which mandate logging security controls within risk management frameworks (RMF). 
Starting in the late 1970s, working groups began establishing criteria for managing auditing and monitoring programs, laying the groundwork for modern cybersecurity practices, such as insider threat detection and incident response. A key publication during this period was NIST’s Special Publication 500-19.
In 2005, the term "SIEM" (Security Information and Event Management) was introduced by Gartner analysts Mark Nicolett and Amrit Williams. SIEM systems provide a single interface for gathering security data from information systems and presenting it as actionable intelligence. The National Institute of Standards and Technology provides the following definition of SIEM: "Application that provides the ability to gather security data from information system components and present that data as actionable information via a single interface." In addition, NIST has designed and implemented a federally mandated RMF.
With the implementation of RMFs globally, auditing and monitoring have become central to information assurance and security. Cybersecurity professionals now rely on logging data to perform real-time security functions, driven by governance models that incorporate these processes into analytical tasks. As information assurance matured in the late 1990s and into the 2000s, the need to centralize system logs became apparent. Centralized log management allows for easier oversight and coordination across networked systems.
On May 17, 2021, U.S. President Joseph Biden signed Executive Order 14028, "Improving the Nation's Cybersecurity," which established further logging requirements, including audit logging and endpoint protection, to enhance incident response capabilities. This order was a response to an increase in ransomware attacks targeting critical infrastructure. By reinforcing information assurance controls within RMFs, the order aimed to drive compliance and secure funding for cybersecurity initiatives.


== Information assurance ==
Published in September 2006, the NIST SP 800-92 Guide to Computer Security Log Management serves as a key document within the NIST Risk Management Framework to guide what should be auditable. As indicated by the absence of the term "SIEM", the document was released before the widespread adoption of SIEM technologies. Although the guide is not exhaustive due to rapid changes in technology since its publication, it remains relevant by anticipating industry growth. NIST is not the only source of guidance on regulatory mechanisms for auditing and monitoring, and many organizations are encouraged to adopt SIEM solutions rather than relying solely on host-based checks.
Several regulations and standards reference NIST’s logging guidance, including the Federal Information Security Management Act (FISMA), Gramm-Leach-Bliley Act (GLBA), Health Insurance Portability and Accountability Act (HIPAA), Sarbanes-Oxley Act (SOX) of 2002, Payment Card Industry Data Security Standard (PCI DSS), and ISO 27001. Public and private organizations frequently reference NIST documents in their security policies.
NIST SP 800-53 AU-2 Event Monitoring is a key security control that supports system auditing and ensures continuous monitoring for information assurance and cybersecurity operations. SIEM solutions are typically employed as central tools for these efforts. Federal systems categorized based on their impact on confidentiality, integrity, and availability (CIA) have five specific logging requirements (AU-2 a-e) that must be met. While logging every action is possible, it is generally not recommended due to the volume of logs and the need for actionable security data. AU-2 provides a foundation for organizations to build a logging strategy that aligns with other controls.
NIST SP 800-53 SI-4 System Monitoring outlines the requirements for monitoring systems, including detecting unauthorized access and tracking anomalies, malware, and potential attacks. This security control specifies both the hardware and software requirements for detecting suspicious activities. Similarly, NIST SP 800-53 RA-10 Threat Hunting, added in Revision 5, emphasizes proactive network defense by identifying threats that evade traditional controls. SIEM solutions play a critical role in aggregating security information for threat hunting teams.
Together, AU-2, SI-4, and RA-10 demonstrate how NIST controls integrate into a comprehensive security strategy. These controls, supported by SIEM solutions, help ensure continuous monitoring, risk assessments, and in-depth defense mechanisms across federal and private networks.


== Terminology ==
The acronyms SEM, SIM and SIEM have sometimes been used interchangeably, but generally refer to the different primary focus of products:

Log management: Focus on simple collection and storage of log messages and audit trails.
Security information management (SIM): Long-term storage as well as analysis and reporting of log data.
Security event manager (SEM): Real-time monitoring, correlation of events, notifications and console views.
Security information and event management (SIEM): Combines SIM and SEM and provides real-time analysis of security alerts generated by network hardware and applications.
Managed Security Service: (MSS) or Managed Security Service Provider: (MSSP): The most common managed services appear to evolve around connectivity and bandwidth, network monitoring, security, virtualization, and disaster recovery.
Security as a service (SECaaS): These security services often include authentication, anti-virus, anti-malware/spyware, intrusion detection, penetration testing and security event management, among others.
In practice many products in this area will have a mix of these functions, so there will often be some overlap – and many commercial vendors also promote their own terminology. Oftentimes commercial vendors provide different combinations of these functionalities which tend to improve SIEM overall. Log management alone doesn't provide real-time insights on network security, SEM on its own won't provide complete data for deep threat analysis. When SEM and log management are combined, more information is available for SIEM to monitor.
A key focus is to monitor and help manage user and service privileges, directory services and other system-configuration changes; as well as providing log auditing and review and incident response.


== Capabilities ==
Data aggregation: Log management aggregates data from many sources, including networks, security, servers, databases, applications, providing the ability to consolidate monitored data to help avoid missing crucial events.
Correlation: Looks for common attributes and links events together into meaningful bundles. This technology provides the ability to perform a variety of correlation techniques to integrate different sources, in order to turn data into useful information. Correlation is typically a function of the Security Event Management portion of a full SIEM solution.
Alerting: The automated analysis of correlated events.
Dashboards: Tools can take event data and turn it into informational charts to assist in seeing patterns, or identifying activity that is not forming a standard pattern.
Compliance: Applications can be employed to automate the gathering of compliance data, producing reports that adapt to existing security, governance and auditing processes.
Retention: Employing long-term storage of historical data to facilitate correlation of data over time, and to provide the retention necessary for compliance requirements. The Long term log data retention is critical in forensic investigations as it is unlikely that the discovery of a network breach will be at the time of the breach occurring.
Forensic analysis: The ability to search across logs on different nodes and time periods based on specific criteria. This mitigates having to aggregate log information in your head or having to search through thousands and thousands of logs.


== Components ==

SIEM architectures may vary by vendor; however, generally, essential components comprise the SIEM engine. The essential components of a SIEM are as follows:

A data collector forwards selected audit logs from a host (agent based or host based log streaming into index and aggregation point) 
An ingest and indexing point aggregation point for parsing, correlation, and data normalization
A search node that is used for visualization, queries, reports, and alerts (analysis take place on a search node) 
A basic SIEM infrastructure is depicted in the image to the right.


== Use cases ==
Computer security researcher Chris Kubecka identified the following SIEM use cases, presented at the hacking conference 28C3 (Chaos Communication Congress).

SIEM visibility and anomaly detection could help detect zero-days or polymorphic code. Primarily due to low rates of anti-virus detection against this type of rapidly changing malware.
Parsing, log normalization and categorization can occur automatically, regardless of the type of computer or network device, as long as it can send a log.
Visualization with a SIEM using security events and log failures can aid in pattern detection.
Protocol anomalies that can indicate a misconfiguration or a security issue can be identified with a SIEM using pattern detection, alerting, baseline and dashboards.
SIEMS can detect covert, malicious communications and encrypted channels.
Cyberwarfare can be detected by SIEMs with accuracy, discovering both attackers and victims.
Modern SIEM platforms support not only detection, but response too. The response can be manual or automated including AI based response.


== Regulatory compliance ==
SIEM systems play a significant role in helping organizations meet regulatory compliance requirements by centralizing log collection, enabling real-time monitoring, and generating audit-ready reports. Under the HIPAA Security Rule, covered entities must implement audit controls that record and examine activity in information systems containing electronic protected health information (ePHI). A proposed update to the HIPAA Security Rule (NPRM, December 2024) would require covered entities to deploy technology capable of detecting and responding to security incidents within 24 hours, a capability commonly provided by SIEM platforms.
The Payment Card Industry Data Security Standard (PCI DSS) Requirement 10 mandates that organizations track and monitor all access to network resources and cardholder data, with SIEM systems commonly used to aggregate and analyze the required log data across distributed environments. Similarly, the Sarbanes-Oxley Act (SOX) requires organizations to maintain audit trails for financial system access, and SIEM platforms provide the centralized logging and retention capabilities necessary to demonstrate compliance during audits.


== Correlation rules examples ==

SIEM systems can have hundreds and thousands of correlation rules. Some of these are simple, and some are more complex. Once a correlation rule is triggered the system can take appropriate steps to mitigate a cyber attack. Usually, this includes sending a notification to a user and then possibly limiting or even shutting down the system.


=== Brute Force Detection ===
Brute force detection is relatively straightforward. Brute forcing relates to continually trying to guess a variable. It most commonly refers to someone trying to constantly guess your password - either manually or with a tool. However, it can refer to trying to guess URLs or important file locations on your system.
An automated brute force is easy to detect as someone trying to enter their password 60 times in a minute is impossible.

An intrusion detection system (IDS) is a device or software application that monitors a network or systems for malicious activity or policy violations. Any intrusion activity or violation is typically either reported to an administrator or collected centrally using a security information and event management (SIEM) system. A SIEM system combines outputs from multiple sources and uses alarm filtering techniques to distinguish malicious activity from false alarms.
IDS types range in scope from single computers to large networks. The most common classifications are network intrusion detection systems (NIDS) and host-based intrusion detection systems (HIDS). A system that monitors important operating system files is an example of an HIDS, while a system that analyzes incoming network traffic is an example of an NIDS. It is also possible to classify IDS by detection approach. The most well-known variants are signature-based detection (recognizing bad patterns, such as exploitation attempts) and anomaly-based detection (detecting deviations from a model of "good" traffic, which often relies on machine learning). Another common variant is reputation-based detection (recognizing the potential threat according to the reputation scores). Some IDS products have the ability to respond to detected intrusions. Systems with response capabilities are typically referred to as an intrusion prevention system (IPS). Intrusion detection systems can also serve specific purposes by augmenting them with custom tools, such as using a honeypot to attract and characterize malicious traffic.


== Comparison with firewalls ==
Although they both relate to network security, an IDS differs from a firewall in that a conventional network firewall (distinct from a next-generation firewall) uses a static set of rules to permit or deny network connections. It implicitly prevents intrusions, assuming an appropriate set of rules have been defined. Essentially, firewalls limit access between networks to prevent intrusion and do not signal an attack from inside the network. An IDS describes a suspected intrusion once it has taken place and signals an alarm. An IDS also watches for attacks that originate from within a system. This is traditionally achieved by examining network communications, identifying heuristics and patterns (often known as signatures) of common computer attacks, and taking action to alert operators.  A system that terminates connections is called an intrusion prevention system, and performs access control like an application layer firewall.


== Intrusion detection category ==
IDS can be classified by where detection takes place (network or host) or the detection method that is employed (signature or anomaly-based).


=== Analyzed activity ===


==== Network intrusion detection systems ====
Network intrusion detection systems (NIDS) are placed at a strategic point or points within the network to monitor traffic to and from all devices on the network. It performs an analysis of passing traffic on the entire subnet, and matches the traffic that is passed on the subnets to the library of known attacks. Once an attack is identified, or abnormal behavior is sensed, the alert can be sent to the administrator. NIDS function to safeguard every device and the entire network from unauthorized access.
An example of an NIDS would be installing it on the subnet where firewalls are located in order to see if someone is trying to break into the firewall. Ideally one would scan all inbound and outbound traffic, however doing so might create a bottleneck that would impair the overall speed of the network. OPNET and NetSim are commonly used tools for simulating network intrusion detection systems. NID Systems are also capable of comparing signatures for similar packets to link and drop harmful detected packets which have a signature matching the records in the NIDS. When we classify the design of the NIDS according to the system interactivity property, there are two types: on-line and off-line NIDS, often referred to as inline and tap mode, respectively. On-line NIDS deals with the network in real time. It analyses the Ethernet packets and applies some rules, to decide if it is an attack or not. Off-line NIDS deals with stored data and passes it through some processes to decide if it is an attack or not.
NIDS can be also combined with other technologies to increase detection and prediction rates. Artificial Neural Network (ANN) based IDS are capable of analyzing huge volumes of data due to the hidden layers and non-linear modeling, however this process requires time due its complex structure. This allows IDS to more efficiently recognize intrusion patterns. Neural networks assist IDS in predicting attacks by learning from mistakes; ANN based IDS help develop an early warning system, based on two layers. The first layer accepts single values, while the second layer takes the first's layers output as input; the cycle repeats and allows the system to automatically recognize new unforeseen patterns in the network. This system can average 99.9\% detection and classification rate, based on research results of 24 network attacks, divided in four categories: DOS, Probe, Remote-to-Local, and user-to-root.


==== Host intrusion detection systems ====

Host intrusion detection systems (HIDS) run on individual hosts or devices on the network. A HIDS monitors the inbound and outbound packets from the device only and will alert the user or administrator if suspicious activity is detected. It takes a snapshot of existing system files and matches it to the previous snapshot. If the critical system files were modified or deleted, an alert is sent to the administrator to investigate. An example of HIDS usage can be seen on mission critical machines, which are not expected to change their configurations.


=== Detection method ===


==== Signature-based ====
Signature-based IDS is the detection of attacks by looking for specific patterns, such as byte sequences in network traffic, or known malicious instruction sequences used by malware. This terminology originates from anti-virus software, which refers to these detected patterns as signatures. Although signature-based IDS can easily detect known attacks, it is difficult to detect new attacks, for which no pattern is available.
In signature-based IDS, the signatures are released by a vendor for all its products. On-time updating of the IDS with the signature is a key aspect.


==== Anomaly-based ====
Anomaly-based intrusion detection systems were primarily introduced to detect unknown attacks, in part due to the rapid development of malware. The basic approach is to use machine learning to create a model of trustworthy activity, and then compare new behavior against this model. Since these models can be trained according to the applications and hardware configurations, machine learning based method has a better generalized property in comparison to traditional signature-based IDS. Although this approach enables the detection of previously unknown attacks, it may suffer from false positives: previously unknown legitimate activity may also be classified as malicious. Most of the existing IDSs suffer from the time-consuming during detection process that degrades the performance of IDSs. Efficient feature selection algorithm makes the classification process used in detection more reliable.
New types of what could be called anomaly-based intrusion detection systems are being viewed by Gartner as User and Entity Behavior Analytics (UEBA) (an evolution of the user behavior analytics category) and network traffic analysis (NTA). In particular, NTA deals with malicious insiders as well as targeted external attacks that have compromised a user machine or account. Gartner has noted that some organizations have opted for NTA over more traditional IDS.


== Intrusion prevention ==
Some systems may attempt to stop an intrusion attempt but this is neither required nor expected of a monitoring system. Intrusion detection and prevention systems (IDPS) are primarily focused on identifying possible incidents, logging information about them, and reporting attempts.  In addition, organizations use IDPS for other purposes, such as identifying problems with security policies, documenting existing threats and deterring individuals from violating security policies. IDPS have become a necessary addition to the security infrastructure of nearly every organization.
IDPS typically record information related to observed events, notify security administrators of important observed events and produce reports. Many IDPS can also respond to a detected threat by attempting to prevent it from succeeding. They use several response techniques, which involve the IDPS stopping the attack itself, changing the security environment (e.g. reconfiguring a firewall) or changing the attack's content.
Intrusion prevention systems (IPS), also known as intrusion detection and prevention systems (IDPS), are network security appliances that monitor network or system activities for malicious activity. The main functions of intrusion prevention systems are to identify malicious activity, log information about this activity, report it and attempt to block or stop it..
Intrusion prevention systems are considered extensions of intrusion detection systems because they both monitor network traffic and/or system activities for malicious activity. The main differences are, unlike intrusion detection systems, intrusion prevention systems are placed in-line and are able to actively prevent or block intrusions that are detected. IPS can take such actions as sending an alarm, dropping detected malicious packets, resetting a connection or blocking traffic from the offending IP address. An IPS also can correct cyclic redundancy check (CRC) errors, defragment packet streams, mitigate TCP sequencing issues, and clean up unwanted transport and network layer options.


=== Classification ===
Intrusion prevention systems can be classified into four different types:

Network-based intrusion prevention system (NIPS):  monitors the entire network for suspicious traffic by analyzing protocol activity.
Wireless intrusion prevention system (WIPS): monitor a wireless network for suspicious traffic by analyzing wireless networking protocols.
Network behavior analysis (NBA): examines network traffic to identify threats that generate unusual traffic flows, such as distributed denial of service (DDoS) attacks, certain forms of malware and policy violations.
Host-based intrusion prevention system (HIPS): an installed software package which monitors a single host for suspicious activity by analyzing events occurring within that host.


=== Detection methods ===
The majority of intrusion prevention systems utilize one of three detection methods: signature-based, statistical anomaly-based, and stateful protocol analysis.

Signature-based detection: Signature-based IDS monitors packets in the Network and compares with pre-configured and pre-determined attack patterns known as signatures. While it is the simplest and most effective method, it fails to detect unknown attacks and variants of known attacks.
Statistical anomaly-based detection: An IDS which is anomaly-based will monitor network traffic and compare it against an established baseline. The baseline will identify what is "normal" for that network – what sort of bandwidth is generally used and what protocols are used. It may however, raise a False Positive alarm for legitimate use of bandwidth if the baselines are not intelligently configured. Ensemble models that use Matthews correlation co-efficient to identify unauthorized network traffic have obtained 99.73\% accuracy.
Stateful protocol analysis detection: This method identifies deviations of protocol states by comparing observed events with "pre-determined profiles of generally accepted definitions of benign activity". While it is capable of knowing and tracing the protocol states, it requires significant resources.
